\usection{Recomendaciones}
A partir de los resultados obtenidos, se presentan las siguientes recomendaciones orientadas a escalar las capacidades funcionales de la herramienta y a fortalecer su integración efectiva en contextos de aprendizaje reales.

Proveer una interfaz de desarrollo abierta que permita a terceros crear nodos y extensiones personalizadas mediante código nativo. Esto otorgaría una mayor flexibilidad al sistema al permitir trascender la herramienta de subflujos, facilitando además la integración y el soporte para otras placas de hardware populares como Arduino y ESP32, diversificando así las capacidades de experimentación física de los usuarios.

Optimizar la integración de la plataforma con el ecosistema del Internet de las Cosas (IoT), mejorando la capacidad del lenguaje para interactuar con sistemas reactivos y protocolos de comunicación estándar de la industria. Esto permitiría a los estudiantes abordar problemas reales de telemetría y control remoto.

Habilitar capacidades de persistencia de datos mediante la incorporación de funciones para la lectura y escritura de archivos locales. Esto permitiría a los estudiantes interactuar con fuentes de información externas y almacenar los resultados de sus ejecuciones, facilitando la comprensión de conceptos relacionados con el manejo de memoria secundaria y la gestión de estados persistentes dentro de un flujo de datos.

Implementar funcionalidades de trabajo cooperativo en tiempo real que permitan a múltiples usuarios colaborar de manera simultánea sobre un mismo proyecto. Esta característica fortalecería el aprendizaje social y la resolución conjunta de desafíos lógicos, replicando las dinámicas de colaboración síncrona que son esenciales en los entornos de desarrollo de software profesionales modernos.
